% Chinese article typography and scientific report conventions.
\usepackage{iftex}
\RequireXeTeX
\usepackage[a4paper,margin=2.5cm]{geometry}
\usepackage{amsmath,amsthm,mathtools}
\usepackage{graphicx,xcolor}
\usepackage[flushleft]{threeparttable}
\usepackage{booktabs,makecell,siunitx,array,multirow,tabularx}
\usepackage{placeins}
\usepackage{caption}
\usepackage[titles]{tocloft}
\usepackage[
    backend=biber,
    style=gb7714-2015,
    gbnamefmt=lowercase,
    gbpub=false,
    doi=false,
    url=false,
    eprint=false
]{biblatex}
\usepackage[unicode,colorlinks=true,allcolors=blue]{hyperref}

% Explicit font faces keep Chinese weight and mixed-script text consistent.
\setCJKmainfont{FandolSong-Regular.otf}[
    BoldFont=FandolHei-Regular.otf,
    ItalicFont=FandolKai-Regular.otf,
    BoldItalicFont=FandolKai-Regular.otf
]
\setCJKsansfont{FandolHei-Regular.otf}[BoldFont=FandolHei-Regular.otf]
\setCJKmonofont{FandolFang-Regular.otf}
% NewTX initializes its paired Latin text and math families after CJK fonts.
\usepackage{newtx}
\usepackage{bm}
\linespread{1.39}
\setlength{\parskip}{0pt}
\pagestyle{plain}
\sisetup{per-mode=symbol}

\ctexset{
    autoindent=true,
    contentsname={目录},
    abstractname={摘要},
    bibname={参考文献},
    figurename={图},
    tablename={表},
    section={
      format={\normalfont\zihao{-3}\bfseries\raggedright},
      beforeskip={20bp plus 1ex minus 0.2ex},
      afterskip={5bp plus 0.2ex},
      aftername={\quad},
      afterindent=true
     },
    subsection={
            format={\normalfont\zihao{4}\bfseries\raggedright},
            beforeskip={12bp plus 1ex minus 0.2ex},
            afterskip={5bp plus 0.2ex},
            aftername={\quad},
            afterindent=true
        },
    subsubsection={
            format={\normalfont\zihao{-4}\bfseries\raggedright},
            beforeskip={12bp plus 1ex minus 0.2ex},
            afterskip={5bp plus 0.2ex},
            aftername={\quad},
            afterindent=true
        }
}
\setcounter{secnumdepth}{3}
\setcounter{tocdepth}{3}

% The title and abstract form a compact article opening.
\makeatletter
\setlength{\@fptop}{0pt}
\setlength{\@fpsep}{12bp}
\setlength{\@fpbot}{0pt plus 1fil}
\renewcommand{\@maketitle}{%
    \begingroup
    \centering
    {\normalfont\fontsize{25bp}{30bp}\bfseries\linespread{1.15}\selectfont\@title\par}%
    \ifx\@author\@empty\else
        \vspace{8bp}%
        {\normalfont\normalsize\begin{tabular}[t]{c}\@author\end{tabular}\par}%
    \fi
    \ifx\@date\@empty\else
        \vspace{4bp}%
        {\normalfont\small\@date\par}%
    \fi
    \endgroup
    \vspace{12bp}%
}
\renewenvironment{abstract}{%
    \par\begingroup
    {\centering\normalfont\zihao{4}\bfseries\abstractname\par}%
    \vspace{8bp}%
    \normalfont\normalsize
    \setlength{\parindent}{2\ccwd}%
    \ignorespaces
}{%
    \par\endgroup
    \vspace{8bp}%
}
\renewcommand{\tableofcontents}{%
    \begingroup
    \hypersetup{hidelinks}%
    \ctexset{section={%
          format={\normalfont\zihao{-2}\bfseries\centering},%
          beforeskip=0pt,%
          afterskip={18bp plus 0.2ex}%
         }}%
    \pdfbookmark[1]{\contentsname}{report-contents}%
    \section*{\contentsname}%
    \@starttoc{toc}%
    \endgroup
}
\makeatother
\renewcommand*{\bibfont}{\zihao{5}\linespread{1.27}\selectfont}
\setlength{\bibitemsep}{3bp}

% Contents use body-size text, dotted leaders, and character-based nesting.
\setlength{\cftbeforesecskip}{6bp plus 1bp}
\setlength{\cftsubsecindent}{\ccwd}
\setlength{\cftsubsubsecindent}{2\ccwd}
\setlength{\cftsecnumwidth}{1.5em}
\setlength{\cftsubsecnumwidth}{2.3em}
\setlength{\cftsubsubsecnumwidth}{3.2em}
\renewcommand{\cftsecdotsep}{\cftdotsep}

% Captions use the body family and a space after the figure or table number.
\DeclareCaptionFont{reportcaption}{\zihao{5}}
\DeclareCaptionLabelSeparator{quad}{\quad}
\captionsetup{font=reportcaption,labelsep=quad}
\captionsetup[table]{position=top}
\setlength{\belowcaptionskip}{5pt}

\newcommand{\PreserveBackslash}[1]{\let\temp=\\#1\let\\=\temp}
\newcolumntype{C}[1]{>{\PreserveBackslash\centering}p{#1}}
\newcolumntype{R}[1]{>{\PreserveBackslash\raggedleft}p{#1}}
\newcolumntype{L}[1]{>{\PreserveBackslash\raggedright}p{#1}}

\newcommand{\PaperTableStyle}{%
    \centering
    \footnotesize
    \setlength{\tabcolsep}{3pt}%
    \renewcommand{\arraystretch}{1.10}%
}
\newcommand{\PaperTableNotesStyle}{\scriptsize}
\newcommand{\ReportStatus}[1]{\textbf{[#1]}}

% Supplementary items use distinct printed numbers and hyperlink targets.
\newcommand{\beginsupplement}{%
    \clearpage
    \setcounter{section}{0}%
    \renewcommand{\thesection}{S\arabic{section}}%
    \renewcommand{\theHsection}{S.\arabic{section}}%
    \renewcommand{\theHsubsection}{S.\arabic{section}.\arabic{subsection}}%
    \setcounter{equation}{0}%
    \renewcommand{\theequation}{S\arabic{equation}}%
    \renewcommand{\theHequation}{S.\arabic{equation}}%
    \setcounter{figure}{0}%
    \renewcommand{\thefigure}{S\arabic{figure}}%
    \renewcommand{\theHfigure}{S.\arabic{figure}}%
    \setcounter{table}{0}%
    \renewcommand{\thetable}{S\arabic{table}}%
    \renewcommand{\theHtable}{S.\arabic{table}}%
}
